\scanpage{209}
\uline{Corollaire} Soit $\deleted{X}Z$ une partie de $X$. On
suppose (i) $\deleted{\unsure{X}}Z$ loc. constructible dans $X$
(ii) $f:X-S$ séparé loc. de type fini (iii) \deleted{$\unsure{X|Z}\to S$}
\inserted{Pour tt $x\in Z$, et $y'$ gén de $f(x)$, $\exists$ gén $x'\in Z$ de $x$, au dessus de $y'$}
\deleted{est ouvert en les $x'\in\unsure{X|Z}$} (iv) les fibres de
$\deleted{\unsure{X}}Z$ sur $S$ sont \uline{propres}. Alors $\deleted{\unsure{X}}Z-S$ est
qu. cpct.

On pt supposer $S$ affine \unsure{donc|ou} noeth, et on est
ramené à prouver que \inserted{si $S\ne\varnothing$,} il existe un ouvert $U\ne\varnothing$
de $S$ tel que $\deleted{\unsure{X}}Z\cap f^{-1}U$ soit qu. cpct. On peut supposer
\unsure{$S$ irréduct. de pt gén $\eta$.}

On construit un $Z'$ de type fini sur $S$ et un
isom $Z'_\eta\to Z_\eta$, \deleted{on peut supposer $Z'$, propre} \inserted{il existe un ouvert qu. cpct}
$V$ de type fini de $X$ contenant $Z_\eta$. Restreignant $S$,
on peut supposer $Z'_\eta\to Z_\eta\subset V$ provient d'un
morphisme $Z'\to\deleted{\unsure{X|Z}}V$, [que ce dernier est
une immersion], \inserted{que $Z\cap V=Z'$} que $Z'$ est \uline{propre} sur $S$.

Alors $Z'$ est \uline{fermé} dans $X$, et par suite on
aura $Z=(Z\cap V=Z')\cup(Z-Z')$, décomposition
\inserted{de $Z$} en deux \uline{ouverts \inserted{relatifs} disjoints}. On est
ainsi ramené au cas où $Z_\eta=\varnothing$,
qui implique $Z=\varnothing$ ....
% archive sheet number: 208
